% The Faedo–Galerkin method

\section{The Faedo--Galerkin method}\label{p:galerkin}

\status{stub}
\prereqs{energy-estimates, weak-convergence, bochner-spaces}

\subsection*{The idea}

To solve a PDE, solve a finite-dimensional version of it (an ODE system), prove bounds that do not depend on the dimension, and pass to the limit. The a priori estimates are what make this work: the approximations are built so that the same estimates hold for them.

\subsection*{Statement}

\begin{thmbox}{The scheme}
\begin{enumerate}
\item Choose a basis $\{w_k\}$ of the solution space, usually eigenfunctions of the elliptic part of the equation (so the boundary conditions are built in).
\item Look for $u_m(t)=\sum_{k=1}^m c_k(t)w_k$ satisfying the weak form tested against $w_1,\dots,w_m$ only. This is an ODE system for $c_1,\dots,c_m$; solve it locally by ODE theory.
\item Prove the a priori estimates for $u_m$ with constants independent of $m$; they extend the ODE solution to a uniform time interval.
\item Extract a weakly convergent subsequence and show the limit satisfies the equation, using compactness for the nonlinear terms.
\end{enumerate}

For linear parabolic equations this gives existence and uniqueness of weak solutions with the energy regularity.
\end{thmbox}

\subsection*{Why it is plausible}

Testing the approximate equation with $u_m$ itself is allowed, because $u_m$ is a combination of the test functions. So the energy identity holds exactly at the finite level, with the same constants. In nonlinear problems the ODE system has a coefficient matrix (for compressible flow, $(\rho w_i,w_k)$, a weighted mass matrix) that must be invertible; this is one reason to start with density bounded away from zero.

\subsection*{Exercises}

\begin{exercise}
  Galerkin for the heat equation on $(0,\pi)$. With $w_k = \sin(kx)$ and $u_m = \sum_{k\le m}c_k(t)w_k$, write the Galerkin equations for $u_t = u_{xx} + f$, $u = 0$ at the ends. Solve them, and observe that $u_m$ is the truncated Fourier series of the exact solution.
  \begin{hint}
    Testing against $w_k$ decouples: $c_k' = -k^2c_k + f_k$ with $f_k = \frac2\pi\int_0^\pi f\,w_k$. The basis consists of eigenfunctions, so nothing couples different modes.
  \end{hint}
  \begin{solution}
    The weak form is $(u_t, v) + (u_x, v_x) = (f, v)$ for $v \in H^1_0(0,\pi)$. On $(0,\pi)$, $(w_j, w_k) = \frac\pi2\delta_{jk}$ and $(w_j', w_k') = \frac\pi2k^2\delta_{jk}$. With $u_m = \sum_{k\le m}c_kw_k$ and $v = w_j$, $j \le m$:

    \[ \frac\pi2c_j' + \frac\pi2j^2c_j = (f, w_j), \quad\text{i.e.}\quad c_j' = -j^2c_j + f_j,\ \ f_j = \frac2\pi(f,w_j), \quad c_j(0) = \frac2\pi(u_0,w_j). \]

    So $c_j(t) = e^{-j^2t}c_j(0) + \int_0^te^{-j^2(t-s)}f_j(s)\,ds$. Testing the exact equation with $w_j$ shows that the $j$-th sine coefficient of the exact solution solves the same ODE with the same initial value. Hence $u_m$ is the exact solution's Fourier series cut at $m$: an eigenfunction basis decouples the modes completely.
  \end{solution}
\end{exercise}

\begin{exercise}
  Show that the Galerkin equations $(\partial_t u_m, w_k) + (\nabla u_m, \nabla w_k) = (f, w_k)$, $k \le m$, imply the energy identity $\frac12\frac{d}{dt}\norm{u_m}^2 + \norm{\nabla u_m}^2 = (f, u_m)$, with constants independent of $m$.
  \begin{hint}
    Multiply the $k$-th equation by $c_k(t)$ and sum over $k$. This is "testing with $u_m$", allowed because $u_m$ is a combination of the test functions.
  \end{hint}
  \begin{solution}
    Multiply the $k$-th equation by $c_k(t)$ and sum over $k \le m$. By linearity, $(\partial_tu_m, \sum_kc_kw_k) + (\nabla u_m, \nabla\sum_kc_kw_k) = (f, \sum_kc_kw_k)$, i.e. $(\partial_tu_m, u_m) + \norm{\nabla u_m}^2 = (f, u_m)$, and $(\partial_tu_m, u_m) = \frac12\frac{d}{dt}\norm{u_m}^2$. No constant depends on $m$. Every estimate derived from this identity, for instance with Young and Gronwall, is therefore uniform in $m$.
  \end{solution}
\end{exercise}

\begin{exercise}
  Passing to the limit. If $u_m \weakto u$ in $L^2(0,T;H^1_0)$, show $\int_0^T(\nabla u_m,\nabla w)\psi\,dt \to \int_0^T(\nabla u,\nabla w)\psi\,dt$ for fixed $w \in H^1_0$ and $\psi \in C_c^\infty(0,T)$. Explain why $\int_0^T\!\int u_m^2 w\,\psi$ need not converge to the same expression with $u$.
  \begin{hint}
    The map $v \mapsto \int_0^T(\nabla v,\nabla w)\psi$ is a bounded linear functional on $L^2(0,T;H^1_0)$. Squares are not linear: see the $\sin(nx)$ example in \pref{weak-convergence}{weak and weak* convergence}.
  \end{hint}
  \begin{solution}
    $\ell(v) = \int_0^T(\nabla v, \nabla w)\psi\,dt$ is linear and $|\ell(v)| \le \norm{\nabla w}_{L^2}\norm{\psi}_{L^2(0,T)}\norm{v}_{L^2(0,T;H^1_0)}$, so $\ell$ is a bounded linear functional on $L^2(0,T;H^1_0)$ and $\ell(u_m) \to \ell(u)$ by definition of weak convergence.

    The map $v \mapsto \int_0^T\!\!\int v^2w\psi$ is not linear, and weak convergence does not pass through it. Example: fix $v \in H^1_0$, $v \ne 0$, and $u_m(x,t) = \sin(mt)v(x)$. Then $u_m \weakto 0$ in $L^2(0,T;H^1_0)$, but $\int_0^T\!\!\int u_m^2w\psi = \int_0^T\sin^2(mt)\psi(t)\,dt\int v^2w \to \frac12\int_0^T\psi\int v^2w$, which is not $0$ in general.
  \end{solution}
\end{exercise}

\subsection*{Where to find a proof}

\begin{itemize}
\item L. C. Evans, \emph{Partial Differential Equations}, 2nd ed., AMS, 2010, §7.1.2 (Theorems 1--3, linear parabolic equations).
\item M. Padula, Existence of global solutions for 2-dimensional viscous compressible flows, \emph{J. Funct. Anal.} 69 (1986) 1--20.
\item H. J. Choe, H. Kim, Strong solutions of the Navier--Stokes equations for isentropic compressible fluids, \emph{J. Differential Equations} 190 (2003) 504--523.
\end{itemize}

\neededfor
